\subsection{Sprint 3}

Planejamos os épicos Perfil de Usuário, Conexões e Ao Vivo. Meu
\textit{squad} ficou com Conexões, que incluía seguir e bloquear usuários,
enviar pedidos de amizade e outras interações entre pessoas no aplicativo.
A pedido do professor, procurei assumir menos implementação e me concentrar
na gestão e no apoio aos colegas.

Durante a \textit{sprint}, uma estudante de Sistemas de Informação entrou em
contato por intermédio do professor orientador e pediu que experimentássemos
outra forma de retrospectiva. Cada integrante avaliaria os colegas por um
formulário. Nós, de AGES IV, preparamos as perguntas, incentivamos o
preenchimento por todos e distribuímos os retornos individualmente. Queríamos
identificar pontos de melhoria que talvez não aparecessem na conversa em
grupo.

Usamos as aulas para definir datas de conclusão das tarefas e preparar uma
lista das atividades bloqueantes. Com isso, orientávamos as prioridades e
ficávamos disponíveis para ajudar os colegas. A equipe concluiu todo o
escopo, apesar do feriado, e ainda adiantou tarefas da \textit{sprint} final.
Também ajudei os integrantes de AGES III na revisão dos \acp{pr}, devido ao
volume de alterações que chegava.

O principal problema que percebi estava na divisão dos \textit{squads}.
Quando um grupo concluía seu épico, alguns integrantes ficavam ociosos por
um período. Levantamos a hipótese de que atribuir os épicos à equipe inteira
poderia reduzir essa espera. Pretendíamos testar essa organização na etapa
seguinte, mantendo o acompanhamento mais próximo dos grupos.

Na AGES III, eu havia levantado a hipótese de que faltava uma aproximação
com os menos experientes. No Hangy, pude me dedicar mais a esse acompanhamento
ao assumir menos código. Eu estava relutante em seguir a orientação do
professor, mas os colegas avançaram nas tarefas enquanto eu esclarecia
dúvidas e acompanhava os prazos. Fiquei satisfeito por ter confiado mais no
trabalho deles.

Ao final da \textit{sprint}, planejávamos uma última etapa de ajustes, com
área de moderação e outras funcionalidades. Também pretendíamos acrescentar
testes de ponta a ponta e verificações de segurança. Queríamos melhorar a
qualidade do produto e tentar concorrer ao reconhecimento de melhor projeto
da AGES.
